\section{Related Work}

\paragraph{Normal-only industrial anomaly detection.}
MVTec AD established the modern evaluation setting with normal training
images, anomalous test images, and pixel-accurate defect masks
\cite{bergmann2019mvtec}. PaDiM models each spatial patch by a multivariate
Gaussian distribution and uses the Mahalanobis distance for detection
\cite{defard2021padim}. PatchCore stores a coreset of nominal patch features and
performs nearest-neighbor matching, emphasizing high recall on unseen defects
\cite{roth2022patchcore}. DRAEM learns discriminative reconstruction by
synthesizing abnormal regions \cite{zavrtanik2021draem}, while CutPaste creates
self-supervised anomalies through image cuts and pasting
\cite{li2021cutpaste}. These methods show that local representations and
carefully designed normality assumptions are more effective than global
reconstruction alone.

\paragraph{Foundation-model representations.}
DINOv2 learns general-purpose visual features through self-supervised
pretraining and produces patch descriptors that transfer well without
fine-tuning \cite{oquab2024dinov2}. AnomalyDINO exploits these frozen features
for few-shot retrieval and applies PCA-based foreground masking and rotation
augmentation to make a small memory bank more robust
\cite{damm2025anomalydino}. Its principal advantage is training-free operation:
inference requires only a forward pass through the encoder and posterior
nearest-neighbor search.

\paragraph{Vision-language and efficient alternatives.}
AnomalyCLIP aligns image regions with object-agnostic normal and abnormal
prompts, enabling zero-shot transfer \cite{zhou2024anomalyclip}. EfficientAD
targets millisecond-level inspection with a compact student-teacher design
\cite{batzner2024efficientad}. The newer MVTec AD 2 benchmark expands the
problem to harder lighting and contamination scenarios
\cite{heckler2026mvtec2}. We use AnomalyCLIP only as related context because the
required comparison is the template baseline, official AnomalyDINO, and one
improved AnomalyDINO variant.

\paragraph{Datasets and metrics.}
VisA contains twelve industrial categories with high-resolution images and
defect masks \cite{zou2022spotdiff}. Its official 1-class split trains on
normal images only and evaluates both normal and anomalous test images.
Image-level AUROC measures ranking quality independent of threshold. Pixel
AUROC evaluates localization maps over all test images, including normal
images whose masks are entirely zero. Because categories differ in defect
frequency and scale, we average per-category metrics equally rather than
pooling all pixels or images into one global score.
